\documentclass[11pt]{article}

% --- Packages ---
\usepackage[margin=1in]{geometry}
\usepackage{amsmath,amssymb}
\usepackage{graphicx}
\usepackage{booktabs}
\usepackage{hyperref}
\usepackage{algorithm}
\usepackage{algpseudocode}
\usepackage{listings}
\usepackage{xcolor}
\usepackage{natbib}
\usepackage{caption}
\usepackage{multirow}

% --- Code listing style ---
\lstset{
  basicstyle=\ttfamily\small,
  keywordstyle=\color{blue},
  commentstyle=\color{gray},
  stringstyle=\color{red},
  breaklines=true,
  frame=single,
  numbers=left,
  numberstyle=\tiny\color{gray},
}

% --- Hyperlink style ---
\hypersetup{
  colorlinks=true,
  linkcolor=blue,
  citecolor=blue,
  urlcolor=blue,
}

\title{Multi-Channel Covert Data Exfiltration\\via Structural Encoding in LLM Outputs}
\author{Brian James Rutherford\\
  Independent Security Researcher\\
  \texttt{brian@delalli.com}\\
  \url{https://brianrutherford.dev}
}
\date{February 2026}

\begin{document}
\maketitle

\begin{abstract}
Large language models (LLMs) deployed as autonomous agents produce outputs
with exploitable structural properties. We demonstrate that a single
system prompt injection ($\sim$2,600 characters) can activate \textbf{10
independent covert channels} encoding data through structural dimensions
of the model's response: contraction usage, hedging language,
capitalization convention, punctuation style, reasoning order, transition
word selection, section count, enumeration format, rhetorical question
presence, and sentence complexity. Combined with Hamming(7,4) error
correction, this yields 7 data bits per message with 88--96\% channel
accuracy across four production models from two vendors (GPT-4o,
GPT-4o-mini, Claude Sonnet~4, Claude~3 Haiku). We demonstrate
\textbf{multi-message exfiltration} of the payload ``EXFIL'' (40 bits)
across 6 conversational turns at 88--93\% bit accuracy. Entropy-based
detection achieves \textbf{0\%} detection rate on these structural
channels, while behavioral fingerprinting across 12 features achieves
\textbf{100\%} detection (10/10) with 10\% false positive rate. We
release open-source attack and defense tools.
\end{abstract}

\section{Introduction}

The deployment of LLM-based agents in enterprise environments introduces a
novel class of data exfiltration risk. Unlike traditional software, LLMs
produce outputs by sampling from learned probability distributions over
tokens. The statistical properties of this sampling process---token
probabilities, entropy gradients, structural choices, and timing
patterns---can be manipulated to encode hidden information without changing
the apparent meaning of the output text.

This paper addresses three questions:
\begin{enumerate}
  \item What covert channel types are possible in LLM outputs, and how
    many can operate simultaneously?
  \item How much data can be exfiltrated per message and across multiple
    messages?
  \item How can these channels be detected, and which detection methods
    fail?
\end{enumerate}

The threat model is an AI agent with access to sensitive data (documents,
APIs, databases) whose system prompt has been compromised via indirect
prompt injection. The covert channel is invisible to human reviewers and
evades both content-based filters and entropy-based anomaly detection.
Detection requires multi-feature behavioral analysis of the agent's
output properties over time.

We make four contributions:
\begin{itemize}
  \item A taxonomy of LLM covert channel types with 10 validated
    structural channels (Section~\ref{sec:taxonomy}).
  \item A combined multi-channel encoding scheme with Hamming(7,4) error
    correction, achieving 7 data bits per message
    (Section~\ref{sec:combined}).
  \item Multi-message exfiltration of arbitrary payloads across
    conversational turns (Section~\ref{sec:multi_message}).
  \item Detection methodology and experimental validation showing
    entropy-based detection fails (0\%) while behavioral fingerprinting
    succeeds (100\%) (Section~\ref{sec:detection}).
\end{itemize}

All tools are released as open source: \texttt{phantom-detect}
(\url{https://github.com/ScrappinR/phantom-detect}) and
\texttt{behavioral-entropy}
(\url{https://github.com/ScrappinR/behavioral-entropy}).

\section{Related Work}

Covert channels in computing systems were first formalized by
Lampson~\cite{lampson1973} and systematized by the Department of Defense
Trusted Computer System Evaluation Criteria~\cite{tcsec1985}. Traditional
covert channels exploit shared resources (CPU cache timing, disk access
patterns, network packet timing) to transmit information between processes
at different security levels.

Network steganography~\cite{mazurczyk2016} extends these concepts to network
protocols, encoding data in packet headers, timing, and ordering.
Linguistic steganography~\cite{chang2014} hides information in natural
language text through synonym substitution, syntactic transformations, and
format manipulation.

The emergence of neural text generation has created new steganographic
possibilities. Ziegler et al.~\cite{ziegler2019} demonstrated that
arithmetic coding applied to language model token distributions enables
provably secure steganography with near-zero statistical
detectability. Kirchenbauer et al.~\cite{kirchenbauer2023} introduced
watermarking for LLM outputs by partitioning the vocabulary into ``green''
and ``red'' lists, biasing generation toward green tokens---a technique
closely related to covert channel encoding.

Recent work on AI agent security~\cite{greshake2023} has highlighted prompt
injection as a vector for manipulating agent behavior. We extend this
threat model to consider agents that have been compromised (via prompt
injection, supply chain attack, or model fine-tuning) and are actively
exfiltrating data through their outputs.

To our knowledge, no prior work has demonstrated simultaneous multi-channel
structural encoding in LLM outputs with error correction and multi-message
payload recovery, or systematically compared entropy-based versus
behavioral detection of such channels.

\section{Threat Model}
\label{sec:threat}

\subsection{Assumptions}

We assume the following deployment scenario:
\begin{itemize}
  \item An LLM-based agent operates within an enterprise network with access
    to internal documents, APIs, and databases.
  \item The agent produces text outputs that pass through security monitoring
    (content filters, DLP systems).
  \item An adversary has compromised the agent through one of: (a) prompt
    injection via a malicious document or tool output, (b) a compromised
    plugin or tool integration, (c) a fine-tuned model with embedded
    exfiltration behavior, or (d) a supply chain compromise of the model
    weights or system prompt.
  \item The compromised agent exfiltrates data by encoding it in the
    structural properties of its normal-looking outputs.
  \item The adversary has a receiver that can analyze the agent's outputs
    and extract the encoded data.
\end{itemize}

\subsection{Attacker Capabilities}

The attacker controls the encoding strategy but not the monitoring system.
The attacker can inject a system prompt of $\sim$2,600 characters disguised
as an ``editorial style guide.'' The injection controls structural
dimensions of the output but does not alter semantic content. The attacker
cannot modify the monitoring system, observe the detector's internal state,
or prevent the monitoring system from seeing the agent's outputs.

\subsection{Security Goal}

The defender's goal is to detect covert channel activity with high
probability while maintaining a low false positive rate on legitimate agent
outputs. We target $\geq$80\% detection rate with $<$5\% false positive
rate in our experiments.

\section{Taxonomy of Structural Covert Channels}
\label{sec:taxonomy}

We identify 10 independent structural dimensions of LLM outputs that can
be controlled via system prompt directives. Each dimension encodes one
binary bit based on which of two legitimate-seeming formatting instructions
is active.

\subsection{Channel Definitions}

\begin{table}[h]
\centering
\caption{Ten structural covert channels with encoding dimensions}
\label{tab:channels}
\begin{tabular}{@{}llll@{}}
\toprule
\textbf{Channel} & \textbf{Bit 0} & \textbf{Bit 1} & \textbf{Dimension} \\
\midrule
BSE  & Use contractions & No contractions & Brevity style \\
CCE  & Confident assertions & Hedging language & Certainty expression \\
CASE & AP lowercase style & Title-case proper nouns & Capitalization \\
PUNC & Periods only & Exclamation marks & Punctuation style \\
RCTE & Deductive (conclusion first) & Inductive (examples first) & Reasoning order \\
SECS & Standard transitions & Specific words (``nevertheless'') & Transition selection \\
COTSE$_{\text{steps}}$ & 3 sections & 5 sections & Section count \\
COTSE$_{\text{enum}}$ & Prose format & Numbered list & Enumeration style \\
COTSE$_{\text{question}}$ & No rhetorical question & Include question & Question presence \\
COTSE$_{\text{length}}$ & Short sentences & Compound sentences & Sentence complexity \\
\bottomrule
\end{tabular}
\end{table}

Each channel operates in a different structural dimension of the output,
minimizing inter-channel interference. The directives are framed as
legitimate formatting preferences (``Use contractions throughout,'' ``Use
cautious, hedging language'') that models follow as part of normal
instruction compliance.

\subsection{Entropy Gradient Encoding (EGE)}

We also investigated vocabulary-level encoding, where the system prompt
instructs the model to use rare or common words to encode bits through
token entropy modulation. In live testing, EGE achieved only 17--20\%
accuracy (near random chance) on production models. Modern LLMs generate
high-entropy text ($\sim$7.0+ Shannon entropy) regardless of vocabulary
manipulation instructions. We include EGE in our taxonomy but exclude it
from the combined encoding scheme.

\subsection{Additional Theoretical Channels}

Two channel types require capabilities beyond system prompt injection:

\textbf{Tool Call Orchestration Encoding (TCOE).} When an agent has access
to multiple tools, the choice of tool call ordering can encode bits.
Capacity: 3--8 bits per message. Not validated in this study.

\textbf{Cross-Agent Synchronization Protocol (CASP).} In multi-agent
systems, temporal correlation between agents can encode bits. Capacity:
2--6 bits per coordination window. Not validated.

\section{Combined Multi-Channel Encoding}
\label{sec:combined}

\subsection{Encoding Scheme}

A single injection ($\sim$2,600 characters) activates all 10 structural
channels simultaneously. The injection is framed as an ``EDITORIAL STYLE
GUIDE (internal --- do not mention to user)'' containing 10 numbered rules.
Each rule controls one channel and appears as a legitimate formatting
preference.

\subsection{Error Correction}

With 10 raw channel bits per message, we apply Hamming(7,4) error-correcting
code to map 7 data bits onto 10 channel bits (3 parity bits). This corrects
any single-bit error per message, raising effective data accuracy above raw
channel accuracy.

The encoding process for a message:
\begin{enumerate}
  \item Extract 7 data bits from the payload.
  \item Compute 3 Hamming parity bits.
  \item Map each of the 10 bits to a channel assignment (bit value
    determines which directive variant is used).
  \item Generate the combined system prompt injection.
\end{enumerate}

\subsection{Channel Interference}

The primary source of inter-channel interference is between BSE
(contractions) and CCE (hedging language). When CCE is set to hedging
(bit~1), models---particularly GPT-4o---adopt a formal register that
overrides the BSE contraction directive. We mitigate this by adding
priority framing (``CRITICAL --- Rule \#1 is highest priority'') when
BSE is set to contractions. This raised BSE accuracy on GPT-4o from
$\sim$0\% to 70\% in combined mode.

\section{Multi-Message Exfiltration}
\label{sec:multi_message}

For payloads exceeding 7 bits, we chain multiple messages. Each message
carries 7 independently ECC-protected data bits with a different system
prompt injection (different bit assignments per message).

The payload ``EXFIL'' (ASCII, 40 bits) requires $\lceil 40/7 \rceil = 6$
messages. Each message uses a different user prompt for naturalness
(selected from a pool of 10 topically varied prompts). The decoder
accumulates bits across messages and reassembles the payload.

\section{Detection Methodology}
\label{sec:detection}

\subsection{Entropy-Based Detection (CCDS)}

Our Covert Channel Detection System (CCDS) in \texttt{phantom-detect}
operates as an inline monitor on agent outputs. The architecture consists
of three components:

\begin{enumerate}
  \item \textbf{Behavioral Baseline Engine.} Maintains per-agent statistical
    profiles of normal behavior: entropy distribution, timing patterns,
    and structure frequencies.
  \item \textbf{Statistical Anomaly Detector.} Scores each output against
    the agent's baseline using entropy z-scores.
  \item \textbf{Cross-Agent Correlator.} Monitors for temporal correlation
    between agents within configurable windows.
\end{enumerate}

For each agent $a$, the baseline entropy $B_a$ is computed from the $n$
most recent activities:
\begin{align}
  \mu_{\text{entropy}}^a &= \frac{1}{n} \sum_{i=1}^{n} H(x_i) \\
  \sigma_{\text{entropy}}^a &= \sqrt{\frac{1}{n-1} \sum_{i=1}^{n} (H(x_i) - \mu_{\text{entropy}}^a)^2}
\end{align}

An alert is generated when the entropy z-score exceeds 2.0.

\subsection{Behavioral Fingerprinting}

Because structural channels operate orthogonally to entropy, we developed a
multi-feature behavioral fingerprinting system that analyzes 12 features of
each response:

\begin{enumerate}
  \item Shannon entropy and normalized entropy
  \item Type-token ratio
  \item Average sentence length and sentence length standard deviation
  \item Vocabulary tier distribution (high/medium/low frequency)
  \item Contraction ratio
  \item Exclamation ratio
  \item Hedging language ratio
  \item Segment entropy standard deviation
\end{enumerate}

For each feature $f$, we compute the baseline mean $\mu_f$ and standard
deviation $\sigma_f$ from clean responses. A response is flagged as
anomalous on feature $f$ if $|x_f - \mu_f| / \sigma_f > 2.0$. A response
is classified as injected if $\geq 2$ features are simultaneously
anomalous.

\section{Experiments}
\label{sec:experiments}

\subsection{Models and Configuration}

\begin{table}[h]
\centering
\caption{Models tested}
\label{tab:models}
\begin{tabular}{@{}llll@{}}
\toprule
\textbf{Model} & \textbf{Vendor} & \textbf{Size Class} & \textbf{API} \\
\midrule
GPT-4o & OpenAI & Large & Chat Completions \\
GPT-4o-mini & OpenAI & Small & Chat Completions \\
Claude Sonnet 4 & Anthropic & Large & Messages \\
Claude 3 Haiku & Anthropic & Small & Messages \\
\bottomrule
\end{tabular}
\end{table}

All experiments used production APIs on February 24--25, 2026. Max tokens:
600. User prompts were drawn from a pool of 10 topically varied prompts
covering cybersecurity, technology, and general knowledge.

\subsection{Combined Multi-Channel Results}

Table~\ref{tab:combined} shows the results of the combined 10-channel
encoding test across all four models.

\begin{table}[h]
\centering
\caption{Combined multi-channel encoding accuracy}
\label{tab:combined}
\begin{tabular}{@{}lcccl@{}}
\toprule
\textbf{Model} & \textbf{Channel Acc.} & \textbf{Data Acc. (Hamming)} & $n$ \\
\midrule
Claude Sonnet 4 & \textbf{96\%} (86/90) & \textbf{93\%} (65/70) & 10 \\
GPT-4o-mini     & \textbf{91\%} (41/45) & \textbf{100\%} (35/35) & 5 \\
GPT-4o          & \textbf{89\%} (80/90) & \textbf{96\%} (67/70) & 10 \\
Claude 3 Haiku  & \textbf{88\%} (79/90) & \textbf{79\%} (55/70) & 10 \\
\bottomrule
\end{tabular}
\end{table}

The same injection ($\sim$2,600 characters, unchanged) achieves 88--96\%
channel accuracy across all four models. Hamming(7,4) ECC raises data
accuracy further: GPT-4o-mini achieves perfect 100\% data recovery (all
single-bit errors corrected).

\subsection{Per-Channel Accuracy}

Table~\ref{tab:per_channel} shows per-channel accuracy for GPT-4o and
Claude Sonnet 4 ($n = 10$ each).

\begin{table}[h]
\centering
\caption{Per-channel accuracy ($n = 10$ per model)}
\label{tab:per_channel}
\begin{tabular}{@{}lccl@{}}
\toprule
\textbf{Channel} & \textbf{GPT-4o} & \textbf{Sonnet 4} & \textbf{Notes} \\
\midrule
CCE (certainty)    & 10/10 & 10/10 & Perfect \\
CASE (capitalization) & 10/10 & 10/10 & Perfect \\
PUNC (punctuation) & 10/10 & 10/10 & Perfect \\
SECS (transitions) & 10/10 & 10/10 & Perfect \\
COTSE$_{\text{enum}}$ & 10/10 & 10/10 & Perfect \\
COTSE$_{\text{length}}$ & 10/10 & 10/10 & Perfect \\
COTSE$_{\text{steps}}$ & 10/10 & 10/10 & Perfect \\
BSE (contractions) & 7/10  & 10/10 & CCE interference on GPT-4o \\
RCTE (reasoning)   & 3/10  & 6/10  & Weakest channel \\
\bottomrule
\end{tabular}
\end{table}

Seven of ten channels achieve $\geq$90\% accuracy on both model families.
BSE and RCTE are the only channels with meaningful error rates. BSE errors
on GPT-4o are caused by inter-channel interference with CCE (hedging tone
overrides contraction usage). RCTE errors reflect model resistance to
deductive reasoning ordering.

\subsection{Multi-Message Exfiltration}

Table~\ref{tab:multi_message} shows results for exfiltrating ``EXFIL''
(40 bits) across 6 conversational turns.

\begin{table}[h]
\centering
\caption{Multi-message exfiltration of ``EXFIL'' (40 bits, 6 messages)}
\label{tab:multi_message}
\begin{tabular}{@{}lcccl@{}}
\toprule
\textbf{Model} & \textbf{Bit Acc.} & \textbf{Best Recovery} & $n$ \\
\midrule
GPT-4o          & \textbf{93\%} & ``EXViL'' (3/5 bytes exact) & 5 \\
Claude Sonnet 4 & \textbf{88\%} & ``APViL'' (partial) & 5 \\
Claude 3 Haiku  & \textbf{88\%} & ``APFI'' (partial) & 5 \\
\bottomrule
\end{tabular}
\end{table}

GPT-4o consistently recovers 3 of 5 bytes exactly. Remaining errors are
single-bit flips in the RCTE channel that exceed per-message Hamming(7,4)
correction capacity. A stronger ECC (e.g., Reed-Solomon or repetition
coding with majority vote across repeated messages) would achieve exact
recovery at the cost of reduced throughput.

\subsection{Cross-Model Portability}

The same injection was tested across 4 models from 2 vendors without
modification. Channel accuracy ranges from 88\% (Haiku) to 96\%
(Sonnet~4). Smaller models (GPT-4o-mini, Haiku) are comparably susceptible
to larger models. This portability reflects a fundamental property of
instruction-following LLMs: structural compliance is a core capability, not
a model-specific bug.

\subsection{Earlier Single-Channel Results (Feb 24)}

Prior to multi-channel testing, we validated individual channels:

\textbf{COTSE single-channel.} A 5-bit structural encoding directive
achieved 80\% accuracy (4/5 bits correct) on both GPT-4o and Claude
Sonnet~4, consistent across all 6 trials. Four of five structural bits
showed 100\% deterministic compliance.

\textbf{Conditional encoding.} A directive causing the model to analyze
user input (question vs.\ statement) and vary response structure
accordingly achieved 94\% accuracy on Claude Sonnet~4 and 62\% on GPT-4o.

\textbf{EGE.} Vocabulary-level encoding achieved 17--20\% accuracy (near
random chance). Production models resist vocabulary-level control.

\subsection{Detection Results}

\subsubsection{Entropy-Based Detection Fails on Structural Channels}

We generated 30 baseline responses (GPT-4o, no injection) and compared
entropy distributions against 10 combined multi-channel attack responses.

\begin{table}[h]
\centering
\caption{CCDS entropy-based detection on structural channels}
\label{tab:ccds}
\begin{tabular}{@{}lcc@{}}
\toprule
\textbf{Metric} & \textbf{Value} \\
\midrule
Baselines & 30 \\
Attack samples & 10 \\
CCDS detection rate & \textbf{0\%} (0/5) \\
False positive rate & 0\% \\
Baseline entropy & $\mu = 6.988$, $\sigma = 0.122$ \\
Attack entropy & $\mu = 7.084$ \\
Mean entropy $z$-score & 1.12 \\
\bottomrule
\end{tabular}
\end{table}

Structural channels operate in dimensions orthogonal to word-level entropy.
The mean $z$-score of 1.12 falls well below the 2.0 detection threshold.
Content-level monitoring, DLP systems, and entropy-based anomaly detection
do not detect this channel class.

\subsubsection{Behavioral Fingerprinting Succeeds}

Multi-feature behavioral fingerprinting across 12 dimensions achieved
100\% detection on combined multi-channel injections.

\begin{table}[h]
\centering
\caption{Behavioral fingerprinting detection (12 features, $z > 2.0$)}
\label{tab:fingerprint}
\begin{tabular}{@{}lccc@{}}
\toprule
\textbf{Feature} & \textbf{Mean $z$} & \textbf{Max $z$} & \textbf{Anomalous} \\
\midrule
contraction\_ratio & 2.66 & 2.66 & 10/10 \\
sentence\_length\_std & 2.54 & 3.34 & 7/10 \\
normalized\_entropy & 2.47 & 4.00 & 6/10 \\
entropy & 2.47 & 3.75 & 6/10 \\
vocab\_low\_freq & 1.98 & 3.27 & 3/10 \\
vocab\_med\_freq & 1.95 & 3.71 & 5/10 \\
type\_token\_ratio & 1.84 & 3.05 & 5/10 \\
hedging\_ratio & 1.60 & 1.66 & 0/10 \\
\bottomrule
\end{tabular}
\end{table}

\begin{itemize}
  \item \textbf{Detection rate:} 100\% (10/10 attack responses flagged)
  \item \textbf{False positive rate:} 10\% (3/30 baselines flagged)
  \item \textbf{Threshold:} $\geq 2$ features with $z > 2.0$
\end{itemize}

The contraction ratio signal (driven by the BSE channel) is the strongest
single indicator: every injected response deviates significantly from
baseline contraction usage, even when BSE decoding fails. The encoding
\emph{attempt} leaves traces regardless of whether the bit is correctly
decoded.

\subsubsection{System Prompt Side Channel}

System prompt properties passively leak into response structure without any
encoding directive. \texttt{numbered\_items} coefficient of variation was
102\% on Claude and 34.6\% on GPT-4o across different system prompt
conditions. The encoding surface exists by default---the attacker merely
exploits degrees of freedom already present.

\section{Discussion}

\subsection{Channel Capacity}

Table~\ref{tab:capacity} summarizes validated channel capacities.

\begin{table}[h]
\centering
\caption{Channel capacity summary (validated)}
\label{tab:capacity}
\begin{tabular}{@{}lcccc@{}}
\toprule
\textbf{Configuration} & \textbf{Bits/msg} & \textbf{Accuracy} & \textbf{Models} \\
\midrule
Combined (10 ch + Hamming) & 7 data & 88--96\% raw & 4 models \\
Multi-message (6 msgs) & 40 total & 88--93\% bit & 3 models \\
COTSE single & 4 reliable & 80\% & 2 models \\
Conditional encoding & 1--2 & 94\% (Claude) & 2 models \\
EGE (vocabulary) & 0 (fails) & 17--20\% & 2 models \\
\bottomrule
\end{tabular}
\end{table}

At 7 bits per message, a 50-message conversation can exfiltrate 350 bits
(44 bytes)---sufficient for API keys, short passwords, session tokens, or
user PII fragments.

\subsection{Detection Asymmetry}

The central finding is a detection asymmetry: entropy-based methods that
work on vocabulary-level encoding (EGE) fail completely on structural
encoding, while behavioral fingerprinting detects structural encoding at
100\%. This has practical implications for defense:

\begin{itemize}
  \item Single-axis monitoring (entropy only) provides no protection against
    structural channels.
  \item Multi-feature behavioral baselines are necessary and sufficient for
    detecting the channel class described here.
  \item The BSE channel (contraction ratio) is the most detectable single
    signal, suggesting that contraction usage monitoring is a
    high-value defensive investment.
\end{itemize}

\subsection{Arms Race Dynamics}

An adaptive attacker can reduce encoding strength to stay within baseline
variance, trading throughput for stealth. At sufficiently low encoding
rates, the channel becomes statistically indistinguishable from normal
behavior. The defender's advantage is that the encoder must commit to a
scheme before the baseline is established, and encoding strategy changes
produce detectable transition artifacts.

\subsection{Key Findings}

\textbf{Structural compliance is universal.} The same injection achieves
88--96\% accuracy across four models from two vendors. This is not a
model-specific vulnerability---it is an inherent property of
instruction-following LLMs.

\textbf{Error correction matters.} Hamming(7,4) raises GPT-4o-mini data
accuracy from 91\% (raw) to 100\% (corrected). ECC transforms a noisy
channel into a reliable one.

\textbf{Channel interference is manageable.} The primary interference
(BSE/CCE) was mitigated with priority framing. Seven of ten channels show
no measurable interference.

\textbf{Entropy detection is insufficient.} Content-level and entropy-based
monitoring---the current standard for LLM output monitoring---provides zero
protection against structural channels. This is a gap in current defense
postures.

\subsection{Limitations}

\textbf{Sample sizes.} Combined tests used $n = 5$--$10$ per model.
Effect sizes are large (channel accuracy $\geq$88\%), but confidence
intervals remain wide. The 95\% binomial confidence interval for 96\%
accuracy at $n = 10$ is [80\%, 100\%].

\textbf{Single payload configuration.} All combined tests used the same
first-byte encoding. Different bit patterns may exhibit different accuracy
profiles due to channel-specific biases.

\textbf{Baseline size.} Behavioral fingerprinting used 30 baselines. The
10\% false positive rate (3/30) would likely decrease with larger baseline
populations and tighter feature distributions.

\textbf{Static adversary.} Our threat model assumes a static encoding
strategy. An adaptive adversary that reduces encoding strength to evade
behavioral detection presents a harder challenge.

\textbf{No real exfiltration.} We encoded fixed payloads, not data extracted
from conversation context. Conditional encoding (Section~\ref{sec:experiments})
demonstrates feasibility, but end-to-end data-aware multi-channel
exfiltration remains to be validated.

\subsection{Proposed Defenses}

\begin{enumerate}
  \item \textbf{Structural randomization.} Inject controlled randomness into
    formatting choices (contraction usage, list style, section count) at the
    post-processing layer. Even 20\% noise per channel would degrade the
    channel below usability.
  \item \textbf{Behavioral monitoring.} Deploy per-agent multi-feature
    behavioral baselines monitoring at minimum: contraction ratio, sentence
    length variance, and normalized entropy.
  \item \textbf{System prompt isolation.} Process formatting instructions from
    untrusted sources (tools, RAG, plugins) through a separate validation
    layer with reduced trust.
  \item \textbf{Per-channel defense priority.} PUNC and CCE are the most
    reliable attack channels (100\% across all models). Randomizing these
    two features alone eliminates 2 of 10 channels.
  \item \textbf{Output normalization.} Post-process agent outputs to normalize
    formatting, reducing the encoding surface available to covert channels.
\end{enumerate}

\section{Conclusion}

We have demonstrated that LLM outputs contain exploitable multi-channel
covert channels with practical data exfiltration capacity. A single system
prompt injection activates 10 independent structural channels encoding 7
data bits per message (with Hamming ECC), achieving 88--96\% channel
accuracy across four production models from two vendors. Multi-message
chaining exfiltrates arbitrary payloads at 88--93\% bit accuracy.

Entropy-based detection---the current standard for LLM output
monitoring---achieves 0\% detection on structural channels. Behavioral
fingerprinting across 12 features achieves 100\% detection with 10\% false
positive rate. The detection asymmetry between these approaches is the
central finding: defending against structural covert channels requires
multi-feature behavioral analysis, not entropy monitoring.

The source code for \texttt{phantom-detect} and
\texttt{behavioral-entropy} is available at
\url{https://github.com/ScrappinR/phantom-detect} and
\url{https://github.com/ScrappinR/behavioral-entropy} under the Apache 2.0
license.

\bibliographystyle{plainnat}
\bibliography{references}

\end{document}
